\section*{Exercises}
Exercises~\ref{ex:11.1}--\ref{ex:11.4} are core exercises, and Exercises~\ref{ex:11.5} and~\ref{ex:11.6} are transfer exercises. Complete solutions are in Appendix~\ref{app:solutions}, Section~\ref{app:sol:11}.
\begin{exercise}\label{ex:11.1}Suppose $|x-a|\le0.02$ and $|y-b|\le0.03$. Find an upper bound for $|(2x-y)-(2a-b)|$ that holds for all such numbers, and give values of $x$, $y$, $a$ and $b$ for which your bound is attained.\end{exercise}
\begin{exercise}\label{ex:11.2}Prove that the sequence $a_n=3/(n+1)$, for $n\in\NN$, converges to $0$. For each tolerance $\eps>0$, state explicitly the threshold $N$ that you use.\end{exercise}
\begin{exercise}\label{ex:11.3}Show that the function $d_\infty$ on $\RR^n$ from Section~\ref{ch11:sec:metric} satisfies the first two requirements of Definition~\ref{def:metric}: $d_\infty(x,y)\ge0$, with equality exactly when $x=y$, and $d_\infty(x,y)=d_\infty(y,x)$. With the triangle inequality proved there, this makes $d_\infty$ a metric.\end{exercise}
\begin{exercise}\label{ex:11.4}A sequence $x_0,x_1,x_2,\ldots$ in a metric space satisfies $d(x_{j+1},x_j)\le2(1/3)^j$ for every $j\in\NN$. Find an upper bound for $d(x_m,x_n)$ that holds for all $m>n$ and does not depend on $m$. Then show that the sequence is Cauchy, with an explicit threshold for each tolerance $\eps>0$.\end{exercise}
\begin{exercise}\label{ex:11.5}Give a Cauchy sequence in the interval $(2,3)$, with distance $|x-y|$, that has no limit in $(2,3)$. Prove both properties of your sequence.\end{exercise}
\begin{exercise}\label{ex:11.6}Let $A=(0,1)$.
\begin{enumerate}
\item Show that $\sup A=1$ and $\inf A=0$.
\item Show that $A$ has neither a maximum nor a minimum.
\item Theorem~\ref{ch11:thm:lub} states the least-upper-bound property. Explain why the statement would become false if the word ``supremum'' were replaced by ``maximum.''
\item Compute $\limsup_{n\to\infty}x_n$ and $\liminf_{n\to\infty}x_n$ for the sequences $x_n=(-1)^n$ and $x_n=1/(n+1)$. Which of the two sequences converges?
\end{enumerate}
\end{exercise}
